\section{Methods}
\label{sec:methods}

\subsection{Datasets}
\label{sec:methods_data}
\textbf{UCSD Stress} (primary). Resting-state EEG from
\citep{komarov2020stress}: 17 graduate students, 92 recordings with
DASS-21 and Daily Stress Scale (DSS) self-reports per session. We use
the 70 recordings with both DASS-21 and DSS available. Epochs are
5\,s at 200\,Hz $\times$ 30 channels; the per-recording binary label
is a DASS-21 threshold at $50$, matching the Wang et al.\
\citep{wang2025stress} protocol.

\textbf{EEGMAT} \citep{eegmat} (paired within-subject comparator).
36 healthy subjects $\times$ 2 recordings (resting baseline plus
serial-subtraction mental arithmetic); 19 channels, 60\,s segments.
This dataset provides the \emph{strong within-subject
alpha-desynchronisation contrast} side of our paired experiment
(Sec.~\ref{sec:results_paired}).

\textbf{ADFTD} \citep{adftd} (for subject-atlas cross-dataset
coverage). Alzheimer's disease vs healthy controls, 195 recordings /
65 subjects / 19 channels.

\textbf{TDBRAIN} \citep{tdbrain} (for subject-atlas cross-dataset
coverage). MDD vs HC, 734 recordings / 359 subjects / 26 channels
(downsampled to 19 for FM compatibility).

ADFTD is used in this paper for subject-atlas variance
decomposition (Sec.~\ref{sec:results_atlas}) and as a single
suggestive-but-not-controlled cross-subject anchored datapoint
(Sec.~\ref{sec:results_fm_adds}). TDBRAIN is used for subject-atlas
variance decomposition and additionally as the single
intermediate-anchor case study discussed in
Sec.~\ref{sec:discussion}, where its role in refining the binary
anchored/bounded SDL classification is explained.

\subsection{Pipeline and foundation models}
\label{sec:methods_pipeline}
All three FMs are wrapped in a unified frozen-extractor interface
producing fixed-dimensional embeddings from (channels $\times$ time)
windows: LaBraM \citep{labram2024} 200-d, CBraMod
\citep{cbramod2025} 200-d, REVE \citep{reve2025} 512-d. Weights are
loaded from the official pretrained releases. Input normalisation is
model-dependent to match each FM's pretraining distribution:
$z$-score for LaBraM, $\mu$V-scale unchanged for CBraMod and REVE.

\subsection{Cross-validation and evaluation protocols}
\label{sec:methods_cv}
Primary evaluation is \textbf{subject-level StratifiedGroupKFold(5)},
grouping by \texttt{patient\_id}. Trial-level StratifiedKFold(5) is
included only as a literature-comparison reference (e.g., reproducing
Wang et al.\ \citep{wang2025stress}) and reported transparently
alongside subject-level numbers; it carries subject leakage and
inflates reported BA on this dataset by $20$--$30$\,pp
(Sec.~\ref{sec:results_audit}).

All training uses cuDNN determinism. Seed variance on the UCSD Stress
$70$-recording / $14$-positive regime nevertheless remains
$\pm 5$--$10$\,pp single-seed, consistent with recent small-$N$
stability literature \citep{smalln2025biorxiv}. All headline claims
are reported as mean $\pm$ sample standard deviation (divide by
$n-1$) over at least three random seeds.

\textbf{Permutation null.} Recording-level label shuffle prior to CV.
A batched helper runs ten permutation seeds per FM; real FT and null
FT are compared by a one-sided Mann--Whitney test on the pooled
balanced-accuracy distributions.

\textbf{Class imbalance handling for classical baselines.} Random
Forest, Logistic Regression ($L_2$), and SVM all use
\texttt{class\_weight = `balanced'}; gradient-boosted trees were
retroactively fixed with \texttt{sample\_weight} computed via
\texttt{sklearn.utils.compute\_sample\_weight} after an initial
methodology review found that only this one classical baseline had
lacked balancing. Results before and after this correction are
reported side-by-side in Sec.~\ref{sec:results_audit}.

\subsection{Variance decomposition and representational similarity}
\label{sec:methods_variance}
For the subject-atlas quantification
(Sec.~\ref{sec:results_atlas}) we compute:
\textbf{(i) Pooled label fraction}: nested sum-of-squares
decomposition of frozen FM embeddings into label-mean,
subject-within-label, and residual components, reported as
$\mathrm{SS}_\text{label} / \mathrm{SS}_\text{total}$.
\textbf{(ii) Mixed-effects ICC} with subject as random intercept on
frozen embedding principal components.
\textbf{(iii) RSA}: Spearman correlation between the embedding
representational dissimilarity matrix (RDM) and the subject-identity
RDM, and separately the label-identity RDM. We tabulate the sign of
(subject-$r$ $-$ label-$r$) across 3 FMs $\times$ 4 datasets = 12
cells.
\textbf{(iv) PermANOVA} (499 permutations) on the recording-level
embeddings with subject as the factor.
Five auxiliary subject-separability metrics (silhouette, Fisher,
nearest-neighbour subject-recovery accuracy, LogME, H-score) are
reported in the supplementary for robustness.

\subsection{Contrast-strength anchoring protocol}
\label{sec:methods_anchor}
A naïve operationalisation of the SDL thesis would define ``contrast
strength'' in terms of FM fine-tuning outcomes and then use those
outcomes to validate the thesis --- a circular construction. We
avoid this by triangulating contrast strength from three
FM-independent dimensions, assessed per dataset prior to any FM
fine-tuning.

\textbf{Dimension 1 --- Published within-subject correlate.}
Literature search for a published per-recording within-subject EEG
correlate of the target label, with documented spectral locus and
effect size.

\textbf{Dimension 2 --- Empirical cluster-corrected contrast.}
Group-level cluster-permutation $F$-test on log-PSDs (Welch, 2\,s
segments, 50\% overlap, 0.5--45\,Hz, MNE
\texttt{permutation\_cluster\_test}, 1000 permutations, default
threshold at $p = 0.05$, independent adjacency across channels and
frequency bins) on our own cohort, testing whether any cluster
survives correction at $p < 0.05$.

\textbf{Dimension 3 --- Locus concordance.}
If an empirical cluster survives, does it spatially and spectrally
align with the published within-subject correlate from Dimension 1?
An empirical cluster at an unexpected locus (different band,
different topography) is diagnostically weaker than one matching the
published anchor, because it may reflect a cohort-level confound
rather than a canonical label-relevant within-subject contrast.

A dataset is classified as \emph{Strong-anchored} when all three
dimensions agree (published correlate, significant cluster, and
locus match), \emph{Intermediate} when the empirical cluster
survives but either the published correlate is absent or its
spectral locus does not match, and \emph{Bounded} when no
within-subject correlate is documented and our own cluster test
does not survive correction. Fig.~\ref{fig:sdl_spectrum} uses this
rubric to order datasets. Full cluster-permutation protocols and
topographic summaries for all three datasets we anchor quantitatively
(UCSD Stress, EEGMAT, TDBRAIN) are in App.~A.

The rationale is to falsify SDL if an anchor disagrees with the
behavioural pattern, rather than merely confirm it post hoc.

\subsection{DASS-21 psychometric scope}
\label{sec:methods_dass_scope}
DASS-21 \citep{lovibond1995} is a past-week trait/state screener;
DSS is a daily-stress aggregate. Neither was psychometrically
validated as a per-recording single-session classification label.
We reproduce \citep{wang2025stress}'s per-recording DASS-21
thresholding for independent audit, acknowledging that the failure
of FMs on this label is a joint statement about FM capability and
about instrument--temporal-resolution mismatch. Under SDL this is
expected: the label's within-subject neural correlate is not
established.

\subsection{Hyperparameter selection}
\label{sec:methods_hp}
Hyperparameters for each FM's best fine-tuning configuration on UCSD
Stress were selected via a $3 \times 2 \times 3$ grid (3 learning
rates $\times$ 2 encoder-learning-rate scales $\times$ 3 seeds)
independently per model. The selected configurations differ between
models. The per-model best-hyperparameter numbers are presented as
observations, not as controlled comparisons between architectures.
Sec.~\ref{sec:results_ceiling}'s architecture-independence claim
relies on each architecture at its own best hyperparameters falling
inside a common band, not on a cross-architecture ranking.

\subsection{Reproducibility}
All code, training configurations, frozen and fine-tuned checkpoints,
variance-decomposition outputs, permutation-null logs, the pre/post
LaBraM diagnostic, the class-balanced classical audit, and the
contrast-strength anchoring protocol specification will be released
at the project repository upon publication.
